\documentclass[12pt,a4paper]{report}

\usepackage[a4paper,margin=1in]{geometry}
\usepackage[T1]{fontenc}
\usepackage[utf8]{inputenc}
\usepackage{lmodern}
\usepackage{setspace}
\usepackage{microtype}
\usepackage{hyperref}
\usepackage{xurl}
\usepackage{booktabs}
\usepackage{longtable}
\usepackage{array}
\usepackage{graphicx}
\usepackage{enumitem}
\usepackage{listings}
\usepackage{xcolor}
\usepackage{fancyhdr}
\usepackage{ragged2e}
\usepackage{draftwatermark}
\usepackage{needspace}

\setstretch{1.2}
\setlength{\parskip}{4pt}
\setlength{\parindent}{0pt}
\setlength{\headheight}{15pt}
\addtolength{\topmargin}{-3pt}
\setlength{\emergencystretch}{2em}
\setlength{\LTleft}{0pt}
\setlength{\LTright}{0pt}
\setlist{leftmargin=*, itemsep=3pt, topsep=3pt}

\SetWatermarkText{Bornil Mahmud (241-15-934)}
\SetWatermarkScale{0.33}
\SetWatermarkLightness{0.92}

\newcolumntype{L}[1]{>{\RaggedRight\arraybackslash}p{#1}}

\hypersetup{
    colorlinks=true,
    linkcolor=blue,
    urlcolor=blue,
    breaklinks=true,
    pdftitle={MathLang Solver - Project Report},
    pdfauthor={Bornil Mahmud}
}

\pagestyle{fancy}
\fancyhf{}
\fancyhead[L]{MathLang Solver}
\fancyhead[R]{Project Report}
\fancyfoot[C]{\thepage}

\lstdefinestyle{ccode}{
    language=C,
    basicstyle=\ttfamily\scriptsize,
    breaklines=true,
    frame=single,
    columns=fullflexible,
    keepspaces=true,
    showstringspaces=false,
    breakatwhitespace=true
}

\lstdefinestyle{phpcode}{
    language=PHP,
    basicstyle=\ttfamily\scriptsize,
    breaklines=true,
    frame=single,
    columns=fullflexible,
    keepspaces=true,
    showstringspaces=false,
    breakatwhitespace=true
}

\newcommand{\imgspace}[2]{
\begin{center}
\fbox{\parbox[c][#2][c]{0.9\textwidth}{\centering \textit{Insert image: #1}}}
\end{center}
}

\begin{document}

\begin{titlepage}
\centering
{\Large Final Year Project Report \par}
\vspace{0.8cm}
{\huge\bfseries MathLang Solver \par}
\vspace{0.4cm}
{\Large Natural Language Arithmetic Parser and Web Platform \par}
\vspace{1.2cm}

\begin{tabular}{ll}
Project Name: & MathLang Solver \\
Repository: & math-tokenazaition \\
Submitted By: & Bornil Mahmud \\
Email: & bornilprof@gmail.com \\
Supervisor: & \textit{[Add Supervisor Name]} \\
Institution: & \textit{[Add Institution Name]} \\
Date: & 18 April 2026 \\
\end{tabular}

\vfill
\textbf{Abstract}\\
MathLang Solver is a full-stack software project that accepts English-like arithmetic
sentences, tokenizes and parses them using Flex/Bison, evaluates them through an AST-based C
engine, and delivers authenticated browser access with MySQL-backed history and administration.
This report follows a formal academic format and documents project motivation, methodology,
implementation, results, engineering impact, and future scope.
\end{titlepage}

\pagenumbering{roman}

\chapter*{Course \& Program Outcome}
\addcontentsline{toc}{chapter}{Course \& Program Outcome}

This project maps course-level compiler and software engineering outcomes to practical implementation.

\begin{longtable}{L{0.2\textwidth} L{0.34\textwidth} L{0.4\textwidth}}
\toprule
Outcome & Description & Evidence in Project \\
\midrule
CO-1 & Apply language processing concepts & Flex lexer and Bison grammar in \texttt{src/lexer.l} and \texttt{src/parser.y} \\
CO-2 & Implement robust software modules & C solver core with AST, evaluation, and JSON output \\
CO-3 & Integrate full-stack components & PHP web layer invoking native solver executable \\
CO-4 & Manage data persistence securely & MySQL schema and PDO prepared statements \\
PO-1 & Engineering knowledge & Formal parser pipeline and system design \\
PO-2 & Problem analysis & Natural-language arithmetic interpretation challenge \\
PO-3 & Design/development of solutions & End-to-end architecture from input to result history \\
PO-7 & Environment and sustainability awareness & Lightweight local deployment, minimal runtime footprint \\
PO-8 & Ethics and professionalism & Security policy, role-based access, responsible disclosure \\
PO-9 & Project management & Structured module breakdown and script-based reproducibility \\
\bottomrule
\end{longtable}

\tableofcontents
\clearpage
\pagenumbering{arabic}

\chapter{Introduction}

\section{Introduction}
MathLang Solver is designed to bridge natural language and mathematical evaluation.
Instead of requiring symbolic syntax only, users can submit phrase-based arithmetic commands
such as \texttt{sum of 5 and 3 multiplied by 2}. This project demonstrates the practical application of compiler
theory to solve a real-world usability problem: making arithmetic computation accessible to users who think in natural
language rather than mathematical symbols. The system combines formal language processing techniques (lexical analysis,
syntactic parsing, semantic evaluation) with modern web application architecture, delivering an integrated platform
that is both technically sound and user-friendly. The implementation showcases how classic compiler concepts—lexers, parsers,
and abstract syntax trees---can be leveraged to create expressive and powerful user-facing applications.

\section{Motivation}
Many learners and general users understand arithmetic instructions in words, but conventional
calculators and parsers accept only strict symbolic forms. This project aims to improve usability
while preserving deterministic computational behavior. The motivation stems from observing that students, educators, and
non-technical professionals often struggle with symbolic notation, which creates a cognitive barrier to problem-solving.
By enabling natural language input, the system removes this barrier and allows users to focus on the mathematical
concept rather than syntax rules. Additionally, the step-by-step output provides transparency in computation, making it
ideal for educational settings where understanding the process is as important as obtaining the final answer.
This project thus addresses both accessibility and pedagogical needs in mathematical computation.

\section{Objectives}

\begin{enumerate}
\item Build a grammar-driven parser for sentence-style arithmetic that accepts multiple natural language forms.
\item Generate and evaluate AST structures in C with deterministic evaluation semantics and error handling.
\item Provide transparent and explainable output including tokens, AST, expression form, step-by-step evaluation, and final results.
\item Deliver a secure and intuitive web interface with user authentication, role-based access control, and session management.
\item Persist execution history for both regular users and administrators, enabling tracking, auditing, and future reference.
\item Demonstrate successful integration of formal compiler techniques with practical full-stack web development patterns.
\end{enumerate}

\section{Feasibility Study}

\needspace{4\baselineskip}
\textbf{Technical Feasibility:} High. The required toolchain (GCC, Flex, Bison, PHP, MySQL)
is standard and available on Windows through local setup or XAMPP workflows. All components are mature, well-documented,
and widely used in production systems. The learning curve is moderate, and debugging tools are readily available.
The C language, Flex/Bison ecosystem, and PHP/MySQL stack are proven technologies with extensive community support.

\vspace{3pt}
\textbf{Operational Feasibility:} High. Build and deployment scripts automate build/run operations; a status page provides
runtime validation. The build process is reproducible across different machines, and the deployment model is lightweight
enough for development laptops as well as production servers. No special hardware or infrastructure is required.
The modular architecture allows incremental testing and validation at each phase.

\vspace{3pt}
\textbf{Economic Feasibility:} High. The project uses open-source tooling and local execution, eliminating licensing and
cloud infrastructure costs. The development timeline is realistic for a single developer working part-time over a semester.
Maintenance overhead is low due to the self-contained nature of the system and minimal external dependencies.

\section{Gap Analysis}

The following table identifies key gaps in existing arithmetic computation systems and demonstrates how this project addresses each gap:

\begin{longtable}{L{0.32\textwidth} L{0.3\textwidth} L{0.32\textwidth}}
\toprule
Existing Gap & Typical Limitation & Project Response \\
\midrule
Natural-language input support & Symbolic-only calculators & Grammar-based sentence parsing with multiple forms \\
Explainable solution trace & Final-value only output & Token stream, AST view, and step-by-step intermediate results \\
Role-aware access & No access model & User/admin authentication and fine-grained authorization \\
Execution history & Stateless solve process & MySQL-backed persistent history with timestamps and audit trail \\
Cross-platform reliability & Manual environment setup & Automated build scripts for Windows deployment \\
\bottomrule
\end{longtable}

\section{Project Outcome}

\needspace{3\baselineskip}
The completed system includes a native parsing engine built with Flex and Bison, an authenticated web dashboard for
user-facing interactions, admin management features for system oversight, and persistent solve records stored in MySQL.
The project successfully demonstrates how compiler theory concepts---tokenization, parsing, AST construction, and evaluation---translate
into practical, user-friendly applications. Upon completion, the system is ready for deployment in educational settings,
and the codebase serves as a reference implementation for combining formal language processing with modern web development.
The project also documents lessons learned, architectural decisions, and potential extensions for future development.

\chapter{Proposed Methodology/Architecture}

\section{Requirement Analysis \& Design Specification}

\subsection{Overview}
The design uses separation of concerns:

\begin{itemize}
\item Native core handles lexical analysis, grammar parsing, AST construction, and evaluation.
\item Web core handles authentication, request orchestration, presentation, and storage.
\item Database layer handles users, execution logs, and grammar metadata.
\end{itemize}

\subsection{System Design}

The system flow consists of six sequential steps:

\begin{enumerate}
\item Input sentence arrives from the browser form.
\item PHP validates the session and sanitizes the request path.
\item Solver executable runs with escaped arguments.
\item Flex tokenizes the sentence and Bison applies grammar rules.
\item AST is evaluated and result contract is serialized to JSON.
\item PHP renders output and stores event in the history table.
\end{enumerate}

\section{System Design Image Section}

This section is intentionally image-focused so each major stage can be explained with visual
evidence during viva, presentation, and documentation review.

\subsection{High-Level Workflow Diagram}
\imgspace{Input to Output Pipeline (Browser -> PHP -> C Solver -> JSON -> Database)}{2.2in}

\subsection{Parser Core Block Diagram}
\imgspace{Lexer, Parser, AST, Evaluator Internal Design}{2.2in}

\subsection{Web Layer Flow}
\imgspace{Authentication, Solve Request, History Save, Dashboard Render}{2.2in}

\subsection{Database ER Snapshot}
\imgspace{app\_users, history, grammar\_rules Relationship View}{2.2in}

\subsection{UI Demonstration Sequence}
\imgspace{Login Screen, Solver Panel, Output Panel, History Panel}{2.2in}

\subsection{Admin Monitoring View}
\imgspace{Role-Based Admin Access and Recent Execution Logs}{2.2in}

\section{Overall Project Plan}

\begin{longtable}{L{0.24\textwidth} L{0.34\textwidth} L{0.34\textwidth}}
\toprule
Phase & Work Items & Deliverables \\
\midrule
Phase 1 & Language design and grammar planning & Token/operator specification \\
Phase 2 & C engine implementation & AST, evaluator, JSON output \\
Phase 3 & Web integration & PHP solver bridge and UI workflow \\
Phase 4 & Data layer and role control & MySQL schema, auth, admin access \\
Phase 5 & Testing and refinement & Validation cases and status visibility \\
\bottomrule
\end{longtable}

\chapter{Implementation and Results}

\section{Implementation}

\subsection{Lexer Definitions for Tokenization}

Lexer maps lexical units to parser tokens. It recognizes numbers, operators, and control connectors in a deterministic way.
The lexer also records token flow for explainable debugging and error reporting.

\begin{lstlisting}[style=ccode,caption={Tokenization rules from src/lexer.l}]
[0-9]+(\.[0-9]+)?              {
    yylval.number = strtod(yytext, NULL);
    record("NUMBER", yytext);
    return NUMBER;
}
"add"|"plus"                   { record("ADD", yytext); return ADD; }
"subtract"|"minus"            { record("SUBTRACT", yytext); return SUBTRACT; }
"multiply"|"multiplied"|"times" { record("MULTIPLY", yytext); return MULTIPLY; }
"divide"|"divided"            { record("DIVIDE", yytext); return DIVIDE; }
\end{lstlisting}

\textbf{Input Example:}
\begin{lstlisting}[style=ccode,caption={Raw user sentence input}]
sum of 10 and 5 then divide by 3
\end{lstlisting}

\textbf{Token Output Example:}
\begin{lstlisting}[style=ccode,caption={Illustrative token stream emitted by lexer}]
[SUM:"sum"] [OF:"of"] [NUMBER:"10"] [AND:"and"] [NUMBER:"5"]
[THEN:"then"] [DIVIDE:"divide"] [BY:"by"] [NUMBER:"3"]
\end{lstlisting}

\subsection{Symbol and Function Table Structures}

The project maintains runtime solver state through structured records including tokens,
steps, expression, tree, and error fields. This structure ensures consistent data flow from parsing to output.

\begin{lstlisting}[style=ccode,caption={Solver result contract from src/solver.h}]
typedef struct {
    int ok;
    char *input;
    char *expression;
    char *tree;
    char *error;
    char *suggestion;
    double result;
    char **steps;
    size_t step_count;
    char **tokens;
    size_t token_count;
} SolverResult;
\end{lstlisting}

\textbf{Output Contract Example:}
\begin{lstlisting}[style=ccode,caption={Sample JSON response stored in history}]
{
    "ok": 1,
    "input": "sum of 10 and 5 then divide by 3",
    "expression": "((10 + 5) / 3)",
    "result": 5,
    "steps": ["10 + 5 = 15", "15 / 3 = 5"],
    "tokens": ["SUM", "OF", "NUMBER(10)", "AND", "NUMBER(5)", "THEN", "DIVIDE", "BY", "NUMBER(3)"]
}
\end{lstlisting}

\subsection{Token Definitions and Grammar Rules}

Grammar rules transform the token stream into an AST that preserves operation order.
This stage is the core of correctness for sentence-style arithmetic evaluation.

\begin{lstlisting}[style=ccode,caption={Grammar excerpt from src/parser.y}]
%token <number> NUMBER
%token SUM DIFFERENCE PRODUCT QUOTIENT
%token ADD SUBTRACT MULTIPLY DIVIDE
%token OF BETWEEN AND WITH BY THEN

expression:
      NUMBER { $$ = ast_make_number($1); }
    | expression ADD expression { $$ = ast_make_binary(OP_ADD, $1, $3); }
    | expression SUBTRACT expression { $$ = ast_make_binary(OP_SUBTRACT, $1, $3); }
    | expression MULTIPLY expression { $$ = ast_make_binary(OP_MULTIPLY, $1, $3); }
    | expression DIVIDE expression { $$ = ast_make_binary(OP_DIVIDE, $1, $3); }
;
\end{lstlisting}

\textbf{Input to Parse Example:}
\begin{lstlisting}[style=ccode,caption={Expression forms accepted by grammar}]
Input A: 8 add 4 multiply 2
Input B: product of 8 and 4 then divide by 2
\end{lstlisting}

\textbf{Parser Output Example:}
\begin{lstlisting}[style=ccode,caption={Illustrative AST form}]
Input A AST: (8 + (4 * 2))
Input B AST: ((8 * 4) / 2)
\end{lstlisting}

\subsection{Handling Function Parameters}

Input handling supports both argument-based and stdin-based modes, providing flexibility for different deployment scenarios.

\begin{lstlisting}[style=ccode,caption={Input flow from src/main.c}]
if (argc > 1) {
    input = join_arguments(argc, argv, 1);
}
if (input == NULL) {
    input = read_stdin_all();
}
\end{lstlisting}

\textbf{Execution Mode Example:}
\begin{lstlisting}[style=ccode,caption={CLI invocation and stdin fallback behavior}]
Mode 1 (argument):   solver.exe "sum of 7 and 9"
Mode 2 (stdin):      echo sum of 7 and 9 | solver.exe
\end{lstlisting}

\subsection{Handling Variable Declarations and Assignments}

Direct symbolic variable declarations is not implemented in this version. The current scope
is sentence-based arithmetic expression evaluation. This is an intentional design boundary
for maintaining deterministic grammar control and clear project scope.

\subsection{Expression Handling for Integer Operations}

The evaluator supports add, subtract, multiply, and divide operations over numeric nodes. Integer-like
presentation is preserved when values are whole numbers, improving output readability.

\begin{lstlisting}[style=ccode,caption={Evaluation guard from src/evaluator.c}]
if (node->op == OP_DIVIDE && right == 0.0) {
    parser_set_error("Division by zero.", "Use a non-zero denominator.");
    *ok = 0;
    return 0.0;
}
\end{lstlisting}

\textbf{Evaluation Input/Output Example:}
\begin{lstlisting}[style=ccode,caption={Result examples from evaluator stage}]
Input:  "difference between 25 and 7 then multiply by 2"
Output: result = 36
Steps:  ["25 - 7 = 18", "18 * 2 = 36"]

Input:  "10 divide by 0"
Output: error = "Division by zero."
Hint:   "Use a non-zero denominator."
\end{lstlisting}

\subsection{Executing Statements with Iteration and Decrement}

The current project does not implement general programming statements (loops/decrements)
as a language runtime. Instead, chained arithmetic operations are supported through grammar
connectors such as \texttt{then}, enabling multi-step calculations.

\section{Results and Discussion}

The system successfully parses and evaluates supported sentence forms, returns structured
JSON output with comprehensive metadata, and persists solve history for users. The explainability output significantly
improves traceability compared to a plain calculator response, making it ideal for educational applications.

\subsection{Representative End-to-End Output Snapshot}

\begin{lstlisting}[style=phpcode,caption={Web-visible solve output (simplified)}]
Input    : sum of 14 and 6 then divide by 5
Result   : 4
Tree     : ((14 + 6) / 5)
Steps    : 14 + 6 = 20
           20 / 5 = 4
Status   : Success
\end{lstlisting}

\subsection{Additional Successful Computation Examples}

\begin{lstlisting}[style=phpcode,caption={Diverse arithmetic expression examples}]
Example 1:
Input    : multiply 7 by 8 then add 12
Result   : 68
Steps    : 7 * 8 = 56
           56 + 12 = 68
Status   : Success

Example 2:
Input    : product of 3 and 4 and 5
Result   : 60
Steps    : 3 * 4 = 12
           12 * 5 = 60
Status   : Success

Example 3:
Input    : difference between 100 and 25 then divide by 3
Result   : 25
Steps    : 100 - 25 = 75
           75 / 3 = 25
Status   : Success
\end{lstlisting}

\subsection{Error Handling Output Examples}

\begin{lstlisting}[style=phpcode,caption={Error cases with helpful messages}]
Error Case 1:
Input    : 42 divide by 0
Result   : null
Error    : "Division by zero detected."
Suggestion: "Please use a non-zero denominator."
Status   : Failed

Error Case 2:
Input    : sum and 15 multiply
Result   : null
Error    : "Syntax error: incomplete expression."
Suggestion: "Check that all numbers are provided."
Status   : Failed
\end{lstlisting}

\subsection{Database Persistence Example}

Successful solves are persisted to the MySQL history table with the following structure:

\begin{lstlisting}[style=phpcode,caption={Sample history record in database}]
History Record:
user_id       : 3
input_text    : "sum of 14 and 6 then divide by 5"
result_json   : {"ok": 1, "result": 4, "steps": [...]}
timestamp     : 2026-04-18 15:32:47
solved_successfully : 1
\end{lstlisting}

\subsection{Observed Quality Characteristics}

\begin{itemize}
\item \textbf{Correctness:} Grammar and evaluator produce deterministic and mathematically correct arithmetic outcomes.
\item \textbf{Explainability:} Token stream, AST view, and step traces provide full transparency into computation.
\item \textbf{Robustness:} Invalid tokens, syntax errors, and divide-by-zero conditions return safe, helpful error output.
\item \textbf{Usability:} Natural-language input significantly lowers the barrier for non-technical users.
\item \textbf{Performance:} Parsing and evaluation complete in milliseconds for typical expressions.
\item \textbf{Data Persistence:} All computations are reliably stored with proper timestamps and user association.
\end{itemize}

\section{UI Design}

The web interface includes login/registration, dashboard solve input, output panels, history viewing, and admin controls. The user experience prioritizes simplicity: users log in, enter a natural-language arithmetic expression,
and immediately see the result with step-by-step breakdown. Advanced features like viewing solve history and admin access
are available through role-based navigation. The design supports both casual users (students learning arithmetic) and
power users (educators managing a classroom of problem solves).

\chapter{Engineering Standards and Mapping}

\section{Impact on Society, Environment and Sustainability}

\subsection{Impact on Life}

The project improves accessibility of arithmetic computation for users who prefer natural-language
interaction over symbolic syntax. This is particularly valuable for students with dyscalculia, non-native speakers,
and users with limited mathematical background.

\subsection{Impact on Society \& Environment}

Lightweight local deployment supports educational institutions with limited infrastructure. No heavy compute or cloud dependency
is required for normal operation, making the system suitable for schools with limited IT resources and bandwidth.

\subsection{Ethical Aspects}

The project stores user input and output history; therefore, responsible access control, credential management, and secure deployment practices are essential.
Privacy and data security are treated as first-class concerns throughout the implementation.

\subsection{Sustainability Plan}

The system uses open-source components, modular architecture, and script-based reproducibility, which supports long-term maintainability and transferability.
Future developers can easily extend or modify the system without architectural constraints.

\section{Solo Development Approach}

This project was developed entirely by a single developer, requiring strong self-organization and time management.
The modular architecture---with clear separation of concerns between core parser, web layer, database schema, and deployment scripts---was
essential for managing complexity across multiple technical domains (C, PHP, SQL, system scripting). The modular design also ensures
that the codebase remains maintainable and extensible for future developers or team-based enhancements. By working independently on this full-stack project, the developer
gained deep expertise across the entire technology stack, from low-level lexer/parser implementation to high-level web application design,
security best practices, and database administration. This holistic understanding is valuable for future projects requiring full-stack
integration and architectural decision-making.

\chapter{Conclusion}

\section{Summary}

MathLang Solver demonstrates a complete integration of compiler-style parsing with full-stack web delivery. It transforms sentence-style arithmetic into deterministic and explainable results,
with robust authentication and persistent history. The project successfully bridges the gap between formal language theory and practical
real-world applications, proving that academic concepts like lexical analysis and parsing can directly solve problems in education
and accessibility. The implementation showcases careful attention to software engineering principles including modularity, security,
error handling, and user experience design. By working independently on this full-stack project, the developer has demonstrated capability
across multiple technical domains and the ability to deliver a production-ready system that is both technically sound and user-friendly.

The MathLang Solver project serves as a concrete proof that compiler theory is not merely academic---it directly enables creation
of accessible, usable software that improves lives. The system is ready for deployment in educational institutions and demonstrates
clear value in making mathematics more approachable.

\section{Limitations}

\begin{itemize}
\item Scope limited to basic arithmetic operations (add, subtract, multiply, divide) over real numbers.
\item No symbolic variable model or advanced control flow language (loops, conditionals, functions) as this would require
a significantly more complex grammar and runtime.
\item Security hardening can be extended further with additional controls such as CSRF tokens, rate throttling, and
advanced audit logging.
\item Performance optimization for high-volume concurrent users has not been implemented; the current design is suitable
for educational and small-scale deployments.
\item The grammar is fixed at compile time; runtime grammar extension is not supported in the current version.
\end{itemize}

\section{Future Work and Recommendations}

\begin{itemize}
\item \textbf{Extended Grammar:} Support variables, assignment, and algebraic equation solving (e.g., "solve for x when x plus 5 equals 12").
\item \textbf{Richer Functions:} Add trigonometric functions (sin, cos, tan), logarithms, exponents, and factorial operations
with natural language forms (e.g., "sine of 45 degrees").
\item \textbf{Security Enhancements:} Implement CSRF tokens, rate limiting, stronger input validation, and automated security testing.
\item \textbf{API and Integration:} Expose RESTful API endpoints to enable third-party integrations with learning management systems
and mobile applications.
\item \textbf{Performance Scaling:} Optimize the parser for higher throughput, implement caching strategies, and conduct load testing
for production deployment.
\item \textbf{User Interface:} Develop a mobile-responsive interface and a native mobile application for on-the-go arithmetic solving.
\item \textbf{Analytics:} Collect and visualize learning analytics to help educators identify common misconceptions and knowledge gaps.
\end{itemize}

\chapter*{References}
\addcontentsline{toc}{chapter}{References}

\begin{itemize}
\item Project repository: \url{https://github.com/BornilMahmud/math-tokenazaition}
\item Flex lexical analyzer: \url{https://github.com/westes/flex}
\item GNU Bison parser generator: \url{https://github.com/akimd/bison}
\item PHP source code repository: \url{https://github.com/php/php-src}
\item MySQL database server: \url{https://github.com/mysql/mysql-server}
\item Flex documentation: \url{https://www.google.com/search?q=flex+lexer+documentation}
\item GNU Bison documentation: \url{https://www.google.com/search?q=gnu+bison+parser+documentation}
\item PHP PDO prepared statements: \url{https://www.google.com/search?q=php+pdo+prepared+statements+tutorial}
\end{itemize}

\appendix

\chapter{Project Overview Image Spaces}

\section{System Overview}
\imgspace{System Overview Diagram}{2.3in}

\section{Architecture Diagram}
\imgspace{Complete System Architecture}{2.3in}

\section{UI Screens}
\imgspace{Login and Dashboard Interface}{2.0in}
\imgspace{Solver Output and History Panels}{2.0in}

\section{Admin and Status Screens}
\imgspace{Admin Management Interface}{2.0in}
\imgspace{System Status Monitoring}{2.0in}

\end{document}
